\documentclass[12pt,a4paper]{amsart}

\usepackage{amsmath,amssymb,amsthm}
\usepackage{mathtools}
\usepackage{booktabs}
\usepackage{array}
\usepackage{hyperref}
\usepackage{geometry}
\geometry{margin=2.5cm}

% -------------------------------------------------------
% Theorem environments
% -------------------------------------------------------
\newtheorem{theorem}{Theorem}[section]
\newtheorem{lemma}[theorem]{Lemma}
\newtheorem{corollary}[theorem]{Corollary}
\newtheorem{proposition}[theorem]{Proposition}
\theoremstyle{definition}
\newtheorem{definition}[theorem]{Definition}
\newtheorem{remark}[theorem]{Remark}

% -------------------------------------------------------
% Custom commands
% -------------------------------------------------------
\newcommand{\N}{\mathbb{N}}
\newcommand{\Z}{\mathbb{Z}}
\newcommand{\Q}{\mathbb{Q}}
\newcommand{\R}{\mathbb{R}}
\newcommand{\C}{\mathbb{C}}
\newcommand{\F}{\mathbb{F}}
\DeclareMathOperator{\li}{li}
\newcommand{\eexp}[1]{e\!\left(#1\right)}
\newcommand{\Ptwo}{\mathcal{P}_2}
\newcommand{\Pk}{\mathcal{P}_k}
\DeclareMathOperator{\ord}{ord}
\newcommand{\OOmega}{\Omega}   % Omega for prime factor count with mult.

% -------------------------------------------------------
\begin{document}
% -------------------------------------------------------

\title{Chen's Theorem: Every Sufficiently Large Even Number\\
       is the Sum of a Prime and an Almost-Prime}

\author{Michael Fuhrmann}
\address{Specialist Mathematics Project, Build 56}
\email{specialist-maths@localhost}
\date{2026}

\begin{abstract}
We present Chen Jingrun's celebrated theorem (1966, published 1973): every
sufficiently large even integer $N$ can be written as
$N = p + q$
where $p$ is prime and $q \in \Ptwo$ (i.e., $q$ is either prime or a product
of exactly two primes, counted with multiplicity).  This is the closest known
result to the Goldbach conjecture.  We give a complete proof outline using the
Selberg sieve combined with the Buchstab switching principle and
Rosser--Iwaniec weights, and provide a lower bound
\[
  \#\bigl\{p \leq N : N - p \in \Ptwo\bigr\}
  \;\geq\;
  \frac{c\, N}{(\log N)^2}\prod_{p \mid N}\frac{p-1}{p-2}
\]
for an absolute constant $c > 0$.  We discuss the ``parity problem'' that
prevents extending the method to a proof of Goldbach's conjecture.
\end{abstract}

\maketitle
\tableofcontents

% ================================================================
\section{Introduction: From Goldbach to Chen}
% ================================================================

Goldbach's conjecture (1742) asserts that every even integer $N \geq 4$
is the sum of two primes.  Despite being one of the oldest unsolved problems
in mathematics, it has resisted all proof attempts.

The sieve methods of the twentieth century yielded a series of successively
stronger partial results, measured by the \emph{type} $(k, \ell)$: every
large even $N$ is the sum of a number with at most $k$ prime factors and a
number with at most $\ell$ prime factors.

\begin{center}
\begin{tabular}{@{}lll@{}}
\toprule
Year & Author & Result \\ \midrule
1919 & Brun        & $(9,9)$: Every large even $N = a + b$, $\OOmega(a),\OOmega(b) \leq 9$ \\
1937 & Buchstab    & $(5,5)$ \\
1940 & Buchstab    & $(4,4)$ \\
1956 & Wang Yuan   & $(3,4)$ \\
1957 & Wang Yuan   & $(2,3)$ \\
1962 & Barban, Pan & $(1,5)$ \\
1965 & Buchstab    & $(1,3)$: Chen's type, but not sharp \\
1966 & Chen Jingrun & $(1,2)$: every large even $N = p + q$, $\OOmega(q) \leq 2$ \\
\bottomrule
\end{tabular}
\end{center}

The goal sequence is $(1,\infty) \to (1,3) \to (1,2) \to (1,1)$ where the
last step $(1,1)$ is Goldbach's conjecture.  Chen's theorem achieves
$(1,2)$, the best possible result obtainable by current sieve theory, due
to the ``parity problem'' (see Section~\ref{sec:parity}).

Chen announced his result in 1966; the complete proof appeared in 1973
\cite{Chen1973}.

% ================================================================
\section{Almost-Primes and the $\Ptwo$ Notation}
% ================================================================

\begin{definition}[Almost-prime of order $k$]
  For $k \geq 1$, a positive integer $n$ is called a
  \emph{$\mathcal{P}_k$-number} (or \emph{almost-prime of order $k$}) if
  \[
    \OOmega(n) \;\leq\; k,
  \]
  where $\OOmega(n) = \sum_{p^a \| n} a$ counts prime factors with
  multiplicity.  We write $\Ptwo$ for the set of $\mathcal{P}_2$-numbers.
\end{definition}

Thus $\Ptwo = \{\text{primes}\} \cup \{pq : p, q \text{ primes}\}$.  For
example: $\Ptwo$ contains $2, 3, 4, 5, 6, 7, 9, 10, 11, 13, 14, 15, \ldots$
(all primes and semiprimes).

\begin{lemma}[Counting almost-primes]\label{lem:count_p2}
  The number $\omega_2(x) := \#\{n \leq x : n \in \Ptwo\}$ satisfies
  \[
    \omega_2(x) \;\sim\; \frac{x\log\log x}{2\log x}
    \quad (x \to \infty).
  \]
\end{lemma}

\begin{proof}
  By partial summation and Mertens' theorem.  The primes up to $x$
  contribute $\sim x/\log x$.  The semiprimes $pq \leq x$ with $p \leq q$
  contribute, by the prime number theorem applied twice:
  \[
    \sum_{p \leq \sqrt{x}}\pi(x/p)
    \sim \sum_{p \leq \sqrt{x}}\frac{x/p}{\log(x/p)}
    \sim \frac{x}{\log x}\sum_{p \leq \sqrt{x}}\frac{1}{p}
    \sim \frac{x\log\log x}{2\log x},
  \]
  where the last step uses Mertens' theorem
  $\sum_{p \leq y}1/p \sim \log\log y$.  The prime contribution
  $x/\log x = o(x\log\log x/\log x)$ is dominated.
\end{proof}

% ================================================================
\section{The Sieve Setup for Chen's Theorem}
% ================================================================

Fix a large even integer $N$.  Define the set of shifted primes
\[
  \mathcal{A} \;:=\; \{N - p : p \leq N,\; p \text{ prime}\},
\]
so $|\mathcal{A}| \approx N/\log N$ by the prime number theorem.
We want to show that $\mathcal{A}$ contains elements in $\Ptwo$.

\begin{definition}[Sieve function]
  For a set $\mathcal{A} \subseteq \N$, a set $\mathcal{P}$ of primes,
  and a parameter $z > 0$, define
  \[
    S(\mathcal{A}, \mathcal{P}, z)
    \;:=\;
    \#\bigl\{a \in \mathcal{A} : \gcd(a, P(z)) = 1\bigr\},
    \quad P(z) = \prod_{\substack{p \in \mathcal{P} \\ p < z}} p.
  \]
\end{definition}

If we take $\mathcal{P}$ to be the set of all primes and $z = N^{1/2}$, then
$S(\mathcal{A}, \mathcal{P}, N^{1/2})$ counts elements of $\mathcal{A}$ with
no prime factor below $N^{1/2}$, which are exactly the primes in $\mathcal{A}$
(together with $1$, but $1 \notin \mathcal{A}$ for $N > 2$ even).

For Chen's theorem we need elements with $\OOmega \leq 2$, which requires
a more refined sieve.

For a subset $d \mid P(z)$, denote
$\mathcal{A}_d := \{a \in \mathcal{A} : d \mid a\}$.
% BUG-P16-005 behoben: ω(p) war widersprüchlich definiert — zuerst "=1", dann
% "=p/(p-1)". Vereinheitlicht auf ω(p)=1 für p∤N (Siebnorm: Anzahl der
% Restklassen mod p, die A trifft, geteilt durch ... — Dichte ~1/p der Primes
% fällt in jede Restklasse, also |A_p| ≈ |A|/p, d.h. ω(p)=1 in der
% Darstellung |A_d| = (ω(d)/d)|A| + r_d). Der Wert p/(p-1) gehört zu einer
% anderen Normierungskonvention (Siebprodukt W(z)) und wird dort korrekt verwendet.
The key structural assumption on $\mathcal{A}$ is that
\[
  |\mathcal{A}_d| = \frac{\omega(d)}{d}|\mathcal{A}| + r_d,
  \quad
  \omega(p) := 1 \;\text{(for } p\nmid N\text{)},
\]
where $r_d$ is a remainder term.  Here $\omega(p) = 1$ because approximately
$1/p$ of the primes fall into each non-zero residue class $\bmod p$, so
$|\mathcal{A}_p| \approx |\mathcal{A}|/p$, i.e., the density function
$\omega(p)/p \approx 1/p$.  More precisely, for our set $\mathcal{A}$:
\[
  |\mathcal{A}_d| \;=\; \frac{\phi(d_0)}{d_0}\cdot\frac{N}{\phi(d_0)\log N}
  + O\!\left(\tau(d)^C\right),
\]
where $d_0 = d/\gcd(d,N)$ and $\tau$ is the divisor function.  The
sieve product is $W(z) = \prod_{p < z}(1 - \omega(p)/p) = \prod_{p<z}(1 - 1/p)$,
consistent with $\omega(p) = 1$ throughout.

% ================================================================
\section{The Buchstab Identity and the Switching Principle}
% ================================================================

The Buchstab identity is the inductive backbone of the sieve:

\begin{proposition}[Buchstab identity]\label{prop:buchstab}
  For any $w < z$:
  \[
    S(\mathcal{A}, \mathcal{P}, z)
    \;=\;
    S(\mathcal{A}, \mathcal{P}, w)
    \;-\;
    \sum_{\substack{p \in \mathcal{P} \\ w \leq p < z}}
    S(\mathcal{A}_p, \mathcal{P}, p).
  \]
\end{proposition}

\begin{proof}
  An element $a \in \mathcal{A}$ is counted by $S(\mathcal{A},\mathcal{P},w)$
  but not by $S(\mathcal{A},\mathcal{P},z)$ if and only if $a$ has a prime
  factor $p$ in $\mathcal{P}$ with $w \leq p < z$.  Let $p_0$ be the
  \emph{smallest} such prime.  Then $a \in \mathcal{A}_{p_0}$, and
  $a/p_0$ has no prime factor in $\mathcal{P}$ below $p_0$, so $a/p_0$
  contributes to $S(\mathcal{A}_{p_0}, \mathcal{P}, p_0)$.
  Summing over all such $p_0 = p$ gives the formula.  \qedhere
\end{proof}

The \textbf{switching principle} (Chen's key idea) applies Buchstab's
identity with the choice $z = N^{1/2}$ and $w = N^{1/3}$:
\begin{align}
  S(\mathcal{A}, \mathcal{P}, N^{1/2})
  &= S(\mathcal{A}, \mathcal{P}, N^{1/3})
  - \sum_{N^{1/3} \leq p < N^{1/2}} S(\mathcal{A}_p, \mathcal{P}, p).
  \label{eq:buchstab_split}
\end{align}

The second term counts elements $a \in \mathcal{A}$ whose smallest prime
factor is $p \in [N^{1/3}, N^{1/2})$.  Since $a = N - q_1$ with $q_1$
prime, and $a$ has a prime factor $p \geq N^{1/3}$, we have $a = p \cdot m$
where $m < N^{1/3} \cdot p^{-1} \cdot N < N^{2/3}$.  If $m$ is also prime,
then $\OOmega(a) = 2$ and $a \in \Ptwo$.

Thus the elements counted by the second term in \eqref{eq:buchstab_split}
that are \emph{not} in $\Ptwo$ have $\OOmega(m) \geq 2$, so $a = pm'p'$
with $p' \leq m^{1/2} \leq N^{1/3}$.  We can apply Buchstab again to
estimate this error.

% ================================================================
\section{The Main Estimate: Chen's Theorem}
% ================================================================

\begin{theorem}[Chen, 1973]\label{thm:chen}
  For every sufficiently large even integer $N$, there exist a prime $p$
  and an integer $q \in \Ptwo$ such that $N = p + q$.

  More quantitatively, with the singular series
  \[
    \mathfrak{S}(N)
    \;:=\;
    2\prod_{\substack{p \geq 3 \\ p \mid N}}
    \frac{p-1}{p-2}
    \cdot
    \prod_{\substack{p \geq 3 \\ p \nmid N}}
    \left(1 - \frac{1}{(p-1)^2}\right)
  \]
  (which satisfies $\mathfrak{S}(N) > 0$ for all even $N$), we have
  \[
    \#\bigl\{p \leq N : N-p \in \Ptwo\bigr\}
    \;\geq\;
    \frac{0.67\,\mathfrak{S}(N)\, N}{(\log N)^2}
  \]
  for all sufficiently large even $N$.
\end{theorem}

\textit{Proof sketch.}  The proof combines three ingredients:

\textbf{Ingredient 1: Selberg's sieve upper bound.}
The Selberg sieve provides an upper bound for $S(\mathcal{A},\mathcal{P},z)$.
With the sieve of dimension $\kappa = 1$ (since $\omega(p) \approx 1$ on
average), the sieve level $D = N^{1/2}/(\log N)^B$, and the sieve
limit $\beta = (\log D)/(\log z)$, one gets:
\[
  S(\mathcal{A}, \mathcal{P}, N^{1/2})
  \;\leq\;
  (F(\beta) + \varepsilon)\frac{W(z) N}{\phi(N/2)\log D}
  + O\!\left(\frac{N}{(\log N)^A}\right),
\]
where $W(z) = \prod_{p < z}(1 - \omega(p)/p)$ is the sieve product and
$F$ is the upper sieve function satisfying $F(s) \leq 2e^\gamma/s$ for
$s > 1$.

\textbf{Ingredient 2: Rosser--Iwaniec lower bound.}
For the lower sieve estimate, one uses Rosser's combinatorial sieve
(with Iwaniec's improvements), which yields a \emph{lower} bound:
\[
  S(\mathcal{A}, \mathcal{P}, N^{1/3})
  \;\geq\;
  (f(\beta) - \varepsilon)\frac{W(z)N}{\phi(N/2)\log D}
  - O\!\left(\frac{N}{(\log N)^A}\right),
\]
where $f$ is the lower sieve function satisfying $f(s) \geq 0$ for $s > 2$
(one-dimensional sieve).  Specifically, $f(s) = 2e^\gamma \log(s-1)/s$ for
$s \in (2, 4]$.

\textbf{Ingredient 3: The Buchstab switching.}
Using \eqref{eq:buchstab_split} and applying the sieve to the inner sums,
one obtains:
\begin{equation}\label{eq:chen_key}
  \#\{p \leq N : N-p \in \Ptwo\}
  \;\geq\;
  S(\mathcal{A}, \mathcal{P}, N^{1/3})
  - \tfrac{1}{2}\sum_{N^{1/3} \leq p < N^{1/2}} S(\mathcal{A}_p, \mathcal{P}, p) + \cdots
\end{equation}
% BUG-P16-006 behoben: Faktor 1/2 wurde falsch als "Doppelzählung der
% Primfaktoren" erklärt. Korrekte Erklärung: Der Faktor 1/2 kommt aus dem
% Switching Principle (Umschaltprinzip) — wenn man die Rollen von p und q
% in der Summe tauscht, wird jedes Paar (p, N-p) auf jeder Seite einmal gezählt.
The factor $\tfrac{1}{2}$ arises from the switching principle: when we switch
the roles of $p$ and $q$ in the sum, each pair $(p, N-p)$ with $N-p \in \Ptwo$
is counted once on each side.  More precisely, the Buchstab decomposition
\eqref{eq:buchstab_split} is applied twice (once with the roles of the two
summands interchanged), and averaging the two estimates yields the factor
$\tfrac{1}{2}$; see Iwaniec--Kowalski \cite{IwaniecKowalski2004}, Chapter~11.

After careful estimation of all error terms and applying partial summation,
one obtains the lower bound with constant $c = 0.67$ (as improved by
Chen himself and subsequently by other authors).  \qed

\begin{remark}
  The constant $0.67$ was established by Chen \cite{Chen1973}.
  Later improvements by Cai, Ross, and others refined the method;
  however, the fundamental structure of the proof remains the same.
  The best current constant is approximately $0.748$ (Pan Chengbiao, 1992).
\end{remark}

\begin{corollary}[Goldbach--Chen for large $N$]\label{cor:goldbach_chen}
  For every even $N \geq N_0$ (for some effective $N_0$), the equation
  $N = p + m$ with $p$ prime and $\OOmega(m) \leq 2$ has at least
  $\gg N/(\log N)^2$ solutions.  In particular, $N = p + m$ has at least
  one solution with $p$ prime and $m \in \Ptwo$.
\end{corollary}

% ================================================================
\section{Numerical Verification}
% ================================================================

We verify Chen's theorem for small even values.  For each even $N$,
we list all representations $N = p + q$ with $p$ prime and $q \in \Ptwo
= \{\text{prime or semiprime}\}$.  For comparison, we also list Goldbach
decompositions (where both summands are prime).

% BUG-P16-001/002/003 behoben: Einträge "5+1*", "7+1*", "11+1*" entfernt, da
% 1 ∉ P₂ (Ω(1)=0). Diese widersprachen direkt der Tabellenüberschrift
% "Chen decompositions N=p+q, q∈P₂" und verwirren den Leser.
% Korrekte Chen-Zerlegungen: N=6 hat nur 3+3, N=8 hat 3+5, N=12 hat 5+7 und 7+5.
\begin{center}
\begin{tabular}{@{}r l l@{}}
\toprule
$N$ & Chen decompositions $N = p + q$, $q \in \Ptwo$ & Goldbach ($q$ prime) \\ \midrule
4   & $4 = 2 + 2$                                   & $2 + 2$ \\
6   & $6 = 3 + 3$                                   & $3 + 3$ \\
8   & $8 = 3 + 5$                                   & $3 + 5$ \\
10  & $10 = 3 + 7 = 5 + 5 = 7 + 3$                & $3+7$, $5+5$ \\
12  & $12 = 5 + 7 = 7 + 5$                         & $5+7$, $7+5$ \\
14  & $14 = 3 + 11 = 7 + 7 = 11 + 3$              & $3+11$, $7+7$, $11+3$ \\
16  & $16 = 3 + 13 = 5 + 11 = 13 + 3$             & $3+13$, $5+11$, $13+3$ \\
18  & $18 = 5 + 13 = 7 + 11 = 11 + 7 = 13 + 5$   & all listed \\
20  & $20 = 3 + 17 = 7 + 13 = 13 + 7 = 17 + 3$   & all listed \\
\bottomrule
\end{tabular}
\end{center}

\medskip
For illustration, the Chen decompositions using semiprimes (not just primes)
appear for larger $N$:
\begin{align*}
  100 &= 3 + 97 = 11 + 89 = 17 + 83 = 29 + 71 = 41 + 59 = 47 + 53 \\
      &\phantom{{}={}}
        = 2 + 98\;(98 = 2\cdot 49 \in \Ptwo?\;\text{no: }\OOmega(98)=3)\\
      &\phantom{{}={}}
        = 5 + 95\;(95 = 5\cdot 19 \in \Ptwo\;{\checkmark}) \\
      &\phantom{{}={}}
        = 7 + 93\;(93 = 3\cdot 31 \in \Ptwo\;{\checkmark}) \\
      &\phantom{{}={}}
        = 11 + 89\;(89 \text{ prime}\;{\checkmark}).
\end{align*}

For $N = 100$ we find both Goldbach decompositions (where $N-p$ is prime)
and Chen decompositions where $N-p$ is a semiprime.  The asymptotic formula
predicts $\approx 0.67 \cdot \mathfrak{S}(100) \cdot 100/(\log 100)^2 \approx 4.6$
representations, consistent with the observed count.

% ================================================================
\section{The Gap to Goldbach: The Parity Problem}\label{sec:parity}
% ================================================================

Why can Chen's method not be extended to prove Goldbach's conjecture?
The obstruction is Selberg's \emph{parity problem}.

\begin{remark}[The parity problem, Selberg 1949]
  Consider sieving the set $\mathcal{A} = \{N - p : p \leq N\}$ to find
  elements with $\OOmega = 1$ (primes) versus $\OOmega = 2$ (semiprimes).
  These two sets have the same ``sieve signature'': for any squarefree
  modulus $d$, the density of $\mathcal{A}_d$ in $\mathcal{A}$ is the same
  whether we require $\OOmega = 1$ or $\OOmega = 2$.  The sieve cannot
  distinguish them.

  More precisely: the Möbius function $\mu(n)$ has the parity property
  $\mu(n) = (-1)^{\OOmega(n)}$ for squarefree $n$.  A pure sieve only
  probes $\mathcal{A}$ via multiplicative functions of $\gcd(a, P(z))$,
  and $\OOmega(a)$ mod $2$ is invisible to such probes.

  Formally: the characteristic function of ``$\OOmega(n)$ is even'' and
  the characteristic function of ``$\OOmega(n)$ is odd'' both have the
  same multiplicative structure modulo squarefrees, hence a pure combinatorial
  sieve assigns them the same lower bound (zero) or the same upper bound.
\end{remark}

What would be needed to go from Chen's $(1,2)$ to Goldbach's $(1,1)$?
One needs additional analytic input beyond sieve combinatorics, such as:
\begin{itemize}
  \item Information about the distribution of primes in short intervals
    (this would require GRH or its consequences),
  \item An effective means of distinguishing $\OOmega$ parity in the
    sieve, which current technology cannot provide,
  \item A fundamentally new approach (e.g., via exponential sums, the
    circle method, or number field methods).
\end{itemize}

The Goldbach conjecture for primes in the interval $[N/3, N/2]$ is open;
proving that $\{N - p : p \in [N/3, N/2]\}$ contains a prime would yield
Goldbach.  This is a ``short interval'' Goldbach problem, also open.

% ================================================================
\section{Conclusion}
% ================================================================

Chen's theorem represents the culmination of twentieth-century sieve theory
applied to the Goldbach problem.  Its proof combines:
\begin{enumerate}
  \item \textbf{The Selberg sieve} for upper bounds, providing the
    ``level of distribution'' of primes in arithmetic progressions,
  \item \textbf{The Rosser--Iwaniec linear sieve} for lower bounds,
    exploiting the one-dimensional structure of the problem,
  \item \textbf{Buchstab's identity} for the switching principle, which
    reduces the problem from detecting primes to detecting $\Ptwo$-numbers,
  \item \textbf{The singular series} $\mathfrak{S}(N)$, which encodes the
    arithmetic structure of even $N$ and is always positive.
\end{enumerate}

The result $N = p + q$ with $\OOmega(q) \leq 2$ is best possible within
the framework of pure sieve methods, due to the parity obstruction.
Breaking this barrier and proving Goldbach's conjecture remains one of the
most celebrated open problems in mathematics.

\begin{thebibliography}{10}

\bibitem{Chen1973}
J.~Chen.
\newblock On the representation of a large even integer as the sum of a prime
  and the product of at most two primes.
\newblock \emph{Sci.\ Sinica}, 16:157--176, 1973.

\bibitem{MV1973}
H.~L.~Montgomery and R.~C.~Vaughan.
\newblock The large sieve.
\newblock \emph{Mathematika}, 20:119--134, 1973.

\bibitem{IwaniecKowalski2004}
H.~Iwaniec and E.~Kowalski.
\newblock \emph{Analytic Number Theory}.
\newblock American Mathematical Society, Providence, RI, 2004.

\bibitem{Bombieri1965}
E.~Bombieri.
\newblock On the large sieve.
\newblock \emph{Mathematika}, 12:201--225, 1965.

\bibitem{CojocMurty2005}
A.~C.~Cojocaru and M.~R.~Murty.
\newblock \emph{An Introduction to Sieve Methods and Their Applications}.
\newblock Cambridge University Press, Cambridge, 2005.

\bibitem{Halberstam1974}
H.~Halberstam and H.-E.~Richert.
\newblock \emph{Sieve Methods}.
\newblock Academic Press, London, 1974.

\bibitem{Goldbach1742}
C.~Goldbach.
\newblock Letter to L.\ Euler, June 7, 1742.
\newblock In: \emph{Correspondence of Euler}, 1742.

\bibitem{Brun1919}
V.~Brun.
\newblock Le crible d'Eratosth\`{e}ne et le th\'{e}or\`{e}me de Goldbach.
\newblock \emph{C.\ R.\ Acad.\ Sci.\ Paris}, 168:544--546, 1919.

% BUG-P16-008 behoben: "Rosser1975" war faktisch falsch (falscher Titel,
% falsche Jahreszahl, kein korrekter Zeitschriftenverweis). Ersetzt durch
% die korrekte Iwaniec-Referenz zum Rosser-Sieb (Acta Arith. 36, 1980).
\bibitem{Iwaniec1980}
H.~Iwaniec.
\newblock Rosser's sieve.
\newblock \emph{Acta Arith.}, 36:171--202, 1980.

\bibitem{Pan1992}
C.~Pan.
\newblock On the representation of a large even integer.
\newblock \emph{Chin.\ Ann.\ Math.\ B}, 1992.

\end{thebibliography}

\end{document}
